
% Load the kaobook class
\documentclass[
	fontsize=10pt, % Base font size
	twoside=true, % Use different layouts for even and odd pages (in particular, if twoside=true, the margin column will be always on the outside)
	%open=any, % If twoside=true, uncomment this to force new chapters to start on any page, not only on right (odd) pages
%%	secnumdepth=1, % How deep to number headings. Defaults to 1 (sections)
]{kaobook}

\newcommand{\discipline}{Sciences Industrielles de l'Ingénieur -- PSI $\star$}
\newcommand{\auteur}{Xavier Pessoles}
\newcommand{\repRel}{../../../..}

\input{\repRel/Style/packages_v2}
\newcommand{\repStyle}{\repRel/Style}
\input{\repRel/Style/macros_SII_v2}
\input{\repRel/Style/environment_v2}
\usepackage{\repStyle/UPSTI_pedagogique}


% Choose the language
\usepackage[english,french]{babel}
\frenchsetup{StandardItemLabels=true}



% Load the bibliography package
\usepackage{kaobiblio}
%\addbibresource{Cy_01_ModelisationSystemes.bib} % Bibliography file

% Load mathematical packages for theorems and related environments
\usepackage{kaotheorems}

% Load the package for hyperreferences
\usepackage{kaorefs}

\graphicspath{{images/}{./}} % Paths where images are looked for

\makeindex[columns=3, title=Alphabetical Index, intoc] % Make LaTeX produce the files required to compile the index


\begin{document}

%%----------------------------------------------------------------------------------------
%%	BOOK INFORMATION
%%----------------------------------------------------------------------------------------

\titlehead{Document Template}
\title[SII en PSI* - Xavier Pessoles]{SII en PSI* - Xavier Pessoles}
\author[XP]{Xavier Pessoles}
\date{\today}
%\publishers{An Awesome Publisher}

%----------------------------------------------------------------------------------------
%
\frontmatter % Denotes the start of the pre-document content, uses roman numerals


%%----------------------------------------------------------------------------------------
%%	MAIN BODY
%%----------------------------------------------------------------------------------------
%
\mainmatter % Denotes the start of the main document content, resets page numbering and uses arabic 

\setchapterstyle{kao}

\setcounter{margintocdepth}{\sectiontocdepth}
\marginlayout
\graphicspath{{\repStyle/png}}

\pagestyle{xp.scrheadings}



%%%%%
\newcommand{\repExo}{}
\newcommand{\nomExo}{}
\AtBeginEnvironment{corrige}{\small}

\livrettrue % Livrettrue :  eleve sans les corrections
\collefalse\colletrue
\normaltrue
%%%%%

\newcommand{\exer}{\subsection*}

\proftrue

%%%%%%%%%%%%%%%%%%%%%%%%%%%%%%%%%%%%%%%%%
\renewcommand{\repExo}{\repRel/PSI_Cy_03_ConceptionCommande/Chapitre_01_Correction//}
\renewcommand{\nomExo}{Cy_03_01_TD_PI_03_SysReeduc}

\graphicspath{{\repRel/Style/png/}{\repExo/\nomExo/images/}}

\chapter*{Les interros qui ne comptent pas 01--A}

\question{Tracer le diagramme de Bode de la fonction de transfert $F(p)=\dfrac{10}{p\left( 1+0,3 p\right)}$.}
\ifprof
\begin{center}
\begin{tikzpicture}[
   semilog lines/.style={thin,black!60},
semilog lines 2/.style={thin,black!30},
    xscale=12/5
]

\begin{scope}[yscale=6/300]
\OrdBode{20}
\semilog{-2}{2}{-60}{80}
\BodeGraph[asymp lines,cyan,samples=100] {-2:2}{\POAmpAsymp{10}{0.3}+ \IntAmp{1}}
\BodeGraph{-2:2} {\POAmp{10}{0.3}+ \IntAmp{1}}
\end{scope}

\begin{scope}[yshift=-1.9cm,yscale=1/70]
\UniteDegre
\OrdBode{45}
\semilog{-2}{2}{-180}{0}
\BodeGraph[asymp lines, cyan, samples=100, const plot]{-2:2}{\POArgAsymp{10}{0.3} + \IntArg{1}}
\BodeGraph{-2:2}{\POArg{10}{0.3}+ \IntArg{1}}
\end{scope}
\end{tikzpicture}
\end{center}
\else
\begin{center}
\begin{tikzpicture}[
   semilog lines/.style={thin,black!60},
semilog lines 2/.style={thin,black!30},
    xscale=12/5
]

\begin{scope}[yscale=6/300]
\OrdBode{20}
\semilog{-2}{2}{-60}{80}
%\BodeGraph[asymp lines,cyan,samples=100] {-2:2}{\POAmpAsymp{10}{0.3}+ \IntAmp{1}}
%\BodeGraph{-2:2} {\POAmp{10}{0.3}+ \IntAmp{1}}
\end{scope}

\begin{scope}[yshift=-1.9cm,yscale=1/70]
\UniteDegre
\OrdBode{45}
\semilog{-2}{2}{-180}{0}
%\BodeGraph[asymp lines, cyan, samples=100, const plot]{-2:2}{\POArgAsymp{10}{0.3} + \IntArg{1}}
%\BodeGraph{-2:2}{\POArg{10}{0.3}+ \IntArg{1}}
\end{scope}
\end{tikzpicture}
\end{center}
\fi



\ifprof
\else
\newpage
\fi

\chapter*{Les interros qui ne comptent pas 01--B}

\question{Tracer le diagramme de Bode de la fonction de transfert $F(p)=\dfrac{0,1}{p\left( 1+3 p\right)}$.}
\ifprof
\begin{center}
\begin{tikzpicture}[
   semilog lines/.style={thin,black!60},
semilog lines 2/.style={thin,black!30},
    xscale=12/5
]

\begin{scope}[yscale=6/300]
\OrdBode{20}
\semilog{-2}{2}{-120}{40}
\BodeGraph[asymp lines,cyan,samples=100] {-2:2}{\POAmpAsymp{0.1}{3}+ \IntAmp{1}}
\BodeGraph{-2:2} {\POAmp{0.1}{3}+ \IntAmp{1}}
\end{scope}

\begin{scope}[yshift=-2.8cm,yscale=1/70]
\UniteDegre
\OrdBode{45}
\semilog{-2}{2}{-180}{0}
\BodeGraph[asymp lines, cyan, samples=100, const plot]{-2:2}{\POArgAsymp{0.1}{3} + \IntArg{1}}
\BodeGraph{-2:2}{\POArg{0.1}{3}+ \IntArg{1}}
\end{scope}
\end{tikzpicture}
\end{center}
\else
\begin{center}
\begin{tikzpicture}[
   semilog lines/.style={thin,black!60},
semilog lines 2/.style={thin,black!30},
    xscale=12/5
]

\begin{scope}[yscale=6/300]
\OrdBode{20}
\semilog{-2}{2}{-120}{40}
%\BodeGraph[asymp lines,cyan,samples=100] {-2:2}{\POAmpAsymp{0.1}{3}+ \IntAmp{1}}
%\BodeGraph{-2:2} {\POAmp{0.1}{3}+ \IntAmp{1}}
\end{scope}

\begin{scope}[yshift=-2.8cm,yscale=1/70]
\UniteDegre
\OrdBode{45}
\semilog{-2}{2}{-180}{0}
%\BodeGraph[asymp lines, cyan, samples=100, const plot]{-2:2}{\POArgAsymp{0.1}{3} + \IntArg{1}}
%\BodeGraph{-2:2}{\POArg{0.1}{3}+ \IntArg{1}}
\end{scope}
\end{tikzpicture}
\end{center}
\fi

\end{document}

